% ============================================================
% 预测层重写片段 —— 对齐队友 Cpaper.tex 行文风格
% 按 4 点意见修改：①不讲舍弃只讲选用 ②潜在变量(不只季节)
% ③引用简洁(来源+好处+建模意义) ④补齐队友 TODO
% 替换位置已在各块开头标注（行号对应队友当前 Cpaper.tex）
% ============================================================


% ============================================================
% 【替换块 A】替换原文 行380–444
% 范围：从"预测层需在每日0:00……"起，到光伏模型池(5)LightGBM结束
% 即：预测层引言段 + 两目标时序特性分析 + 异构模型池构建 整段
% ============================================================

预测层需在每日 0:00 时刻，利用截止前一日的全部历史数据，对当日 144 个 10 分钟间隔时点的小区负载功率与光伏发电实际功率进行日前预测；同一天的预测值不回灌增量训练，以避免误差沿 144 个时点逐级累积。

预测的真正难点在于，驱动两条曲线的因素大多不可直接观测。题面只给出历史功率，并未提供气温、辐照、天气类型、节假日等外生变量，但这些潜在变量的影响都沉淀在历史曲线的形态之中：居民作息决定了工作日与周末的周周期差异，季节更替经由气温驱动空调负荷的缓变趋势与峰值时点漂移，阴晴云雨决定光伏逐日出力的短时起伏，日内太阳高度角决定光伏的钟形包络，随机用电与测量误差构成残差噪声。此外，模型还要经历从 31 天冷启动到数百天数据积累的全过程，不同阶段可用信息量差异悬殊。这些因素在不同日期交替主导，没有任何单一模型能对全部情形同时保持最优：统计模型善于刻画稳定周期却难以响应非线性突变，机器学习模型善于融合多因素却依赖充足样本，相似日方法能复现相近工况却受限于历史库规模。本文因此构建机制互补的异构模型池，让各成员分别刻画不同的潜在驱动因素，再由在线权重依据近期表现自动决定各类因素的话语权。

\subsubsection{两目标时序特性分析}

在组合构建前，先分别诊断两条序列的时序特性，为每个成员的"因素分工"提供量化依据。组合若遗漏主导因素，会在对应情境下集体失准；若堆砌机制同质的模型，则类似共线性，徒增算力而无信息增量。

\textbf{小区负载}\ \ 自相关分析显示滞后 7 阶系数达 0.971，而滞后 1 阶仅 0.373，表明周周期是首要成分。周内模式显著分化：周一至周四及周日的日均负荷稳定在 5163--5173~kW，周五、周六骤降至 3264、3267~kW（-36.8\%），对应工作日与周末的居民作息差异。日级方差分解中周季节成分占 89.3\%、趋势占 14.1\%、残差仅 4.4\%。峰值时点还随季节在早、晚高峰间整体迁移：1--5 月与 9--12 月稳定在 09:00--09:50 的早高峰，6--8 月空调季移至 16:40--17:40 的晚高峰，7 月与 11 月的平均日内曲线相差 23.0\%，月度日均从 11 月的 4338~kW 到 7 月的 5337~kW、极值比 1.23；此外节假日明显塌陷，清明、五一、国庆日均较普通工作日低 14.5\%--29.7\%，全年最低的 5 天全部落在节假日及其调休日。

\textbf{光伏发电}\ \ 滞后 1 阶系数为 0.956、滞后 7 阶为 0.878，日连续性远强于周周期；工作日与周末无显著差异，符合其由自然因素驱动的特征。月度日均从 12 月的 1514~kW 升到 8 月的 2916~kW、极值比达 1.93，由季节辐照水平决定。物理约束明确：51.5\% 的时点为夜间零值，全年 66.3\% 的日峰值落在 12:00 前后半小时内，钟形包络稳定。对 365 条白天曲线做 SVD，第一主成分即解释 84.6\% 的方差、前三个主成分累计 94.9\%，说明日曲线共享一条主形态；而剔除季节趋势后，月内日总量的变异系数仅 6.0\%，这部分小幅摆动正是天气噪声。

上述诊断把潜在变量落到了可量化的证据上：负载侧的主导因素是居民作息（周周期占日级方差 89.3\%、周末骤降 36.8\%）、季节气温（峰值时点在早晚高峰间迁移、7 月与 11 月曲线相差 23.0\%）与节假日扰动；光伏侧的主导因素是天文规律（日连续性、SVD 第一主成分占 84.6\%、峰值集中于正午、夜间恒零）与小幅天气噪声（月内日总量变异仅 6.0\%）。两条序列的驱动结构截然不同，故分别构建独立模型池；池内每个成员都有明确分工，分别对某一类潜在因素敏感。

\subsubsection{异构模型池构建}

成员按三条标准选取：机制差异化（各对应一类潜在驱动因素，机制同质的模型误差高度相关，叠加后不带来新信息）、因素分工明确、滚动合规（只使用预测时点之前的数据，杜绝未来信息泄漏）。负载与光伏各设五名成员，下面逐一说明其被选中的理由及所刻画的因素。

\textbf{负载模型池}

（1）\textbf{季节朴素法}：直接以上周同一时点值为预测，专门承载居民作息这一因素。附件中"同一星期几"远比"相邻一天"更可参照（滞后 7 阶自相关 0.971，远高于滞后 1 阶的 0.373），且周一至周四、周日稳定在 5163--5173~kW、周五周六骤降 36.8\%，复制上周曲线即可锁定这套周内规律；它无需训练、不引入参数估计误差，是作息占主导时最稳的强基准，全期 MAE 为 176.80。

（2）\textbf{月度均值预测}：以近 4 周同一时点的平均为预测，用来滤除居民随机用电的残差噪声。附件日级方差中残差仅占 4.4\%，绝大部分波动是可重复的周期成分，多周平均让相互独立的随机扰动彼此抵消、留下稳定的周期期望，全期 MAE 为 188.88。

（3）\textbf{STL 分量预测}：把历史序列分离为趋势与季节后分别外推（趋势阻尼 0.8、季节取上一同相位），承载季节气温驱动的缓变。附件中这类缓变十分明显：峰值时点在 1--5、9--12 月停于早高峰、6--8 月整体迁到晚高峰，7 月与 11 月日内曲线相差 23.0\%、7 月日均较 11 月高 23\%，单靠复制上周无法跟踪，分解出的趋势项却能平滑跟进；其鲁棒拟合同时能压住节假日塌陷这类异常（全年最低 5 天均为节假日调休）~\cite{cleveland1990}，全期 MAE 为 178.59。

（4）\textbf{kNN 相似日}：在近 60 天里挑出与昨日曲线形状最接近的几个同星期几历史日、借用其次日曲线，用来捕捉"同为工作日却彼此不同"的天气与工况因素。这一检索在附件中确实有效：归一化后同星期几候选日与查询的平均曲线距离为 0.269，而最接近的 3 个近邻只有 0.227、贴近 15.7\%，说明曲线形状确能把工况更相似的日子筛出来~\cite{akyurek2014}，全期 MAE 为 188.80。

（5）\textbf{LightGBM}：池中唯一的非线性模型，用 18 维特征把上述因素一并编码——星期几、月份、周末、日内时点等日历特征承载作息与季节，8 个滞后、滑窗特征承载近期趋势与天气惯性，树结构再刻画因素间的交互，例如"空调季 $\times$ 周末 $\times$ 晚高峰"这类周期模型无法表达的组合；附件中峰谷比从 8 月的 1.96 到 12 月的 2.33 逐月变动、周末整体下挫 36.8\%，正是多因素交互的体现。为防小样本过拟合取强正则（num\_leaves=15、min\_child\_samples=30、L1/L2 正则、50 轮早停），它是负载最优单模型，全期 MAE 为 172.71。

\textbf{光伏模型池}

（1）\textbf{Persistence}：直接沿用昨日同刻值，承载由天文规律决定的强日连续性（滞后 1 阶自相关 0.956）。昨日曲线已含当前太阳高度角与近期天气惯性，不做平滑反而能对阴晴突变给出最快响应，全期 MAE 为 185.35。

（2）\textbf{kNN 相似日}：挑出曲线形状最接近的历史日、复现其次日出力，用来匹配天气类型。这一机制在光伏上比负载更显著：最接近 3 邻的归一化曲线距离比全部候选平均低 41.7\%（0.068 对 0.117），曲线形状能清楚区分晴、多云等不同天气型；光伏不受作息支配，故不限定星期几，全期 MAE 为 205.04。

（3）\textbf{STL 分量预测}：分离出辐照水平的季节趋势后外推，承载太阳高度角季节变化主导的缓升缓降——附件月度日均从 12 月 1514~kW 升到 8 月 2916~kW、相差 1.93 倍，正由趋势项刻画，全期 MAE 为 212.73。

（4）\textbf{FPCA 函数型主成分}：对近 14 天白天曲线做低秩重构，以阻尼系数 0.4 融合公共形态与最近一日的变异。附件为这种"一条公共包络加少数变异模态"的结构提供了直接证据：SVD 第一主成分就解释 84.6\% 方差、前三个累计 94.9\%，且 66.3\% 的日峰值集中在 12:00 前后半小时，说明日曲线形状高度共享、天气只是对公共包络做小幅修正（月内日总量变异仅 6.0\%），恰由函数型主成分提取~\cite{wagnermuns2018}，阻尼融合则在保留天气摆动与抑制单日异常间取得平衡，预测值截断为非负。它是光伏最优单模型，全期 MAE 为 160.71。

（5）\textbf{LightGBM}：特征工程与负载侧一致，融合季节辐照水平与近期天气惯性；并利用夜间物理恒零的规律对 06:00 前、20:00 后强制置零（附件中 51.5\% 的时点为夜间零值），全期 MAE 为 192.14。


% ============================================================
% 【替换块 B】替换原文 行449–451（EGD 小节"若仅通过简单平均……"那段）
% 作用：把"季节性"泛化为"多类潜在因素相对重要性随时间交替"
% ============================================================

若仅用简单平均或按全期误差定一组固定权重，便默认各潜在因素的相对重要性全年不变，这与天气、作息、季节交替主导的事实相悖；即使用一个强大的元学习器按全期误差赋权，全期误差也只是各阶段表现的平均，无法对"当前由哪类因素主导"做出差异化响应，反而同时拉低各阶段的预测表现。


% ============================================================
% 【替换块 C】替换原文 行502–505 的 "TODO 表……" 那段
% 即替换 "故据此让初始权重偏向简单方法，具体冷启动结果见下表。 TODO 表负载侧取……"
% ============================================================

而简单平均在高不确定性下反而更稳~\cite{graefe2015}，故让初始权重偏向无需训练的统计成员，具体取值见表~\ref{tab:coldstart}；随着数据积累，式~\eqref{eq:reinit} 的混合系数衰减，机器学习成员的份额逐步交由在线更新决定。

\begin{table}[htbp]
\centering
\caption{冷启动阶段两模型池的初始权重}
{\small
\begin{tabular}{lcc}
\toprule
模型成员 & 负载池 & 光伏池 \\
\midrule
季节朴素法 / Persistence & 0.25 & 0.30 \\
月度均值预测 / kNN       & 0.20 & 0.20 \\
STL 分量预测             & 0.20 & 0.15 \\
kNN / FPCA               & 0.20 & 0.15 \\
LightGBM                 & 0.15 & 0.20 \\
\bottomrule
\end{tabular}}
\label{tab:coldstart}
\end{table}


% ============================================================
% 【删除点 D】删除原文 行568 末尾一句
% 原句："负载池初始 7 模型剔除 FPCA 与谐波回归的精简决策，同样出自这组数据。"
% 删除理由：属于"讲舍弃"，按意见1不提被剔除的模型
% ============================================================


% ============================================================
% 【可选微调 E】原文 行589 "响应月度及季节检验" 段
% 该段本身讲季节响应，保留即可；仅建议把段末"这正是选择在线权重而非固定权重的理由"
% 前补一句点回潜在变量，使论证闭环。可将段末改为：
% ============================================================

% 原段末"……整个切换由式~\eqref{eq:egd} 自动完成，无需人工干预，这正是选择在线权重而非固定权重的理由。"
% 建议改为：
整个切换由式~\eqref{eq:egd} 自动完成，无需人工干预；哪类潜在因素当前占主导，权重就自动向对该因素最敏感的成员倾斜，这正是选择在线权重而非固定权重的根本理由。
